\subsubsection{Deuxième étape: on abandonne les unités dès que c'est trop compliqué}

On note $h_1$ la hauteur du pont par rapport au sol en contrebas de l'élastique. À $h_1=100\un{m}$, lorsque l’homme se laisse tomber, il n’a encore aucune vitesse. Son énergie cinétique est donc nulle : $E_{c_1} = 0\un{J}$ (zéro joule).

Calculons son énergie mécanique (on ne lui donne pas d'index 1 car elle est constante):
\begin{gather*}
    E_m = E_{c_1} + E_{p_1} \\
    E_m = 0\un{J }+ E_{p_1} \\
    E_m = E_{p_1}
\end{gather*}

Supposons que je pèse $70\un{kg}$.
\begin{gather*}
    E_{p_1} = m \times g \times h_1 \\
    \iff E_{p_1} = 70\un{kg} \times 9,81{\frac{N}{\un{kg}}} \times 100\un{m} \\
    \iff E_{p_1} = \SI{68670}{\joule}
\end{gather*}

Lors de la vitesse maximale, 30 mètres plus bas, pour $h_{2} = 100\un{m} - 30\un{m} = 70\un{m}$ de hauteur par rapport au sol:
\begin{gather*}
    E_{p_2} = m \times g \times h_2 \\
    \iff E_{p_2} = 70\un{kg} \times \num{9,81}{\frac{\mathrm{N}}{\mathrm{kg}}} \times 70\un{m} \\
    \iff E_{p_2} = 48069\un{J}
\end{gather*}

Or, l'énergie mécanique est conservée, donc
\begin{gather*}
    E_{c_2} + E_{p_2} = E_{p_1} \\
    \iff E_{c_2} = E_{p_1} - E_{p_2} \\
    \iff E_{c_2} = \SI{68670}{\joule} - \SI{48069}{\joule} = \SI{20601}{\joule}
\end{gather*}
Or $E_{c_2} = \frac{1}{2} \times m \times v^2$, ce nous permet maintenant de chercher $v$:
\begin{gather*}
    \frac{1}{2} \times 70\un{kg} \times v^2 = \SI{20601}{\joule} \\
    \iff v^2 = \frac{20601 \times 2}{70} \\
    \iff v = \sqrt{\num{588.6}} \\
    \boxed{v \approx \SI{24.26}{\meter\per\second}}
\end{gather*}